%% ========================================================================
%% Chapter 1: Introduction to Operating Systems and Linux Installation
%% ========================================================================

\chapter{Introduction to Operating Systems and Linux Installation}
\label{ch:introduction}

\chapterhero{blue}{Understand the foundation of operating systems, Linux distribution choices, the installation process, and initial verification so the system is ready as a learning environment.}{\faDesktop\quad operating system}{installation}{verification}

One program is already active before you can log in, before a browser can be opened, and even before the system asks for a password. That program manages all hardware, shares processor time among many tasks, and ensures that each application receives safe memory space. That program is the \textbf{operating system}. Without an operating system, every program would need to write its own code to communicate with the network card, screen, and disk: an inefficient approach that is almost impossible to apply widely.

This book uses Linux as the learning platform. Linux runs more than 90 percent of internet servers, all of the world's 500 fastest supercomputers, and most global cloud infrastructure. Understanding Linux means understanding the foundation of modern technology.

\textbf{What You'll Learn}

By the end of this chapter, you will be able to:

\begin{enumerate}
    \item explain the definition and main functions of an operating system;
    \item distinguish operating system categories and kernel architectures with examples;
    \item trace the broad evolution of operating systems from the early era to today;
    \item install Ubuntu Server 24.04 LTS in a VirtualBox \term{virtual machine};
    \item explain disk partition concepts and Linux \term{filesystem} types;
    \item perform initial configuration and verify the installation result.
\end{enumerate}

\section{Definition and Functions of an Operating System}
\label{sec:os-basics}

An operating system is software that acts as an intermediary between users and computer hardware. Think of a work desk. On the desk are documents being worked on, writing tools, and space to move. The operating system is the manager of that desk: it decides which documents may be on the desk now, where tools are stored, and who may touch which item.

An operating system has five main functions:

\begin{description}[noitemsep,topsep=2pt,parsep=0pt]
    \item[Process Management] The operating system manages all running programs. It decides when a program receives processor time, creates new processes when needed, and ends processes when they finish. In Linux, \cmd{ps} and \cmd{top} are used to view currently running processes.

    \item[Memory Management] Every active program needs space in RAM. The operating system allocates that space, makes sure one program does not take another program's area, and moves data to disk when RAM starts to fill.

    \item[File Management] The operating system provides a directory structure so data can be stored, named, and found again. It also manages permissions: who may read, modify, or execute a file.

    \item[I/O Management] Devices such as keyboards, printers, and network cards work in very different ways. The operating system provides an abstraction layer through \term{driver} software so programs do not need to understand each device's internal behavior.

    \item[Security and Protection] The operating system verifies user identity during login, limits access to resources based on assigned permissions, and records important activity for security audits.
\end{description}

\begin{notebox}
Without an operating system, every program would need its own \term{driver} for every available hardware device. Imagine rewriting screen control code every time you create a new application. The operating system removes this repeated work.
\end{notebox}

\section{Operating System Types}
\label{sec:os-categories}

Operating systems are developed for many device types with different requirements. Table~\ref{tab:os-categories} summarizes the major categories in use today.

\begin{table}[tbp]
    \centering
    \caption{Operating system categories by platform}
    \label{tab:os-categories}
    \begin{tblr}{
        colspec = {lX[1.6,l]X[1,l]},
        width = \linewidth,
        hlines,
        vlines,
        row{1} = {font=\bfseries},
        cells = {font=\small}
    }
        Category & Characteristics & Examples \\
        Desktop OS &
            Graphical interface for personal computers &
            Windows 11, macOS, Ubuntu Desktop \\
        Server OS &
            Serves many clients; usually without a graphical interface &
            Ubuntu Server, RHEL, Windows Server \\
        Mobile OS &
            Touchscreen and battery power efficiency &
            Android, iOS \\
        Real-Time OS &
            Response within very strict time limits &
            FreeRTOS, QNX, VxWorks \\
        Embedded OS &
            Limited resources; for IoT and routers &
            Embedded Linux, OpenWRT \\
    \end{tblr}
\end{table}

This book focuses on the \textbf{Server OS} category, especially Linux. The skills you gain also apply directly to Linux-based Desktop OS systems and Embedded OS systems that use the same Linux kernel.

\section{Kernel Architecture}
\label{sec:kernel-architecture}

The kernel is the core of an operating system. This program remains in memory from the time the computer starts until it shuts down, and it has direct access to all hardware. Every other part of the operating system, including the user interface and utilities, stays outside the kernel and must request kernel services to access hardware.

\subsection{Kernel Space and User Space}

Modern operating systems divide memory into two strictly separated regions. The relationship between these regions is summarized in \cref{fig:os-layers}. The term \textit{ring} comes from processor privilege levels in CPU architecture. The smaller the number, the higher the privilege. In modern operating system practice, the two most important levels are \textbf{Ring 0} for the kernel and \textbf{Ring 3} for ordinary applications.

\begin{description}[noitemsep,topsep=2pt,parsep=0pt]
    \item[\term{Kernel space} (Ring 0)] The region with the highest privileges. Only kernel code runs here. Code in \term{kernel space} can read and write all memory addresses and communicate directly with hardware.

    \item[\term{User space} (Ring 3)] The region where ordinary applications run. Access to hardware is completely restricted. Applications must pass through a mechanism called a \textbf{system call} to request services from the kernel.
\end{description}

\begin{figure}[tbp]
    \centering
    \resizebox{0.65\textwidth}{!}{%
    \begin{tikzpicture}[node distance=0.8cm]
        \node[draw, rectangle, minimum width=10cm, minimum height=1.2cm,
              fill=blue!10, align=center] (user) {
            \textbf{User Space (Ring 3)}\\[0.2em]
            \small Browser, Text Editor, Media Player, Terminal
        };
        \node[draw, rectangle, minimum width=10cm, minimum height=0.7cm,
              below=0.8cm of user, fill=yellow!20] (syscall) {
            \textbf{System Call Interface}
        };
        \node[draw, rectangle, minimum width=10cm, minimum height=1.6cm,
              below=0.8cm of syscall, fill=red!10, align=center] (kernel) {
            \textbf{Kernel Space (Ring 0)}\\[0.2em]
            \small Process Scheduler, Memory Management, File System, Device Drivers
        };
        \node[draw, rectangle, minimum width=10cm, minimum height=1.1cm,
              below=0.8cm of kernel, fill=gray!20, align=center] (hw) {
            \textbf{Hardware}\\[0.2em]
            \small CPU, RAM, Disk, Network Card, GPU
        };
        \draw[->, thick] (user.south) -- (syscall.north)
            node[midway, right, font=\small] {System Call};
        \draw[->, thick] (syscall.south) -- (kernel.north);
        \draw[->, thick] (kernel.south) -- (hw.north)
            node[midway, right, font=\small] {Direct Access};
    \end{tikzpicture}%
    }
    \caption{Operating system architecture layers}
    \label{fig:os-layers}
\end{figure}

A system call is the only official mechanism applications use to request kernel services. When a program opens a file, for example, it calls the \texttt{open()} system call. The kernel then validates permissions, interacts with the disk \term{driver}, and returns the result to the calling program.

\subsection{Kernel Architecture Types}

There are three main approaches to designing an operating system kernel. Their main characteristics are compared in \cref{tab:kernel-types}, while the placement of service components is shown in \cref{fig:kernel-comparison}.

\begin{table}[tbp]
    \centering
    \caption{Comparison of kernel architecture types}
    \label{tab:kernel-types}
    \begin{tblr}{
        colspec = {lX[l]X[l]X[l]},
        width = \linewidth,
        hlines,
        vlines,
        row{1} = {font=\bfseries},
        cells = {font=\small}
    }
        Aspect & Monolithic & Microkernel & Hybrid \\
        Description &
            All services run in \term{kernel space} inside one large \term{binary} &
            Only essential services stay in the kernel; the rest run in \term{user space} &
            Combines monolithic and microkernel approaches \\
        Performance & Very high & Lower because of communication overhead & High \\
        Reliability &
            One bug can damage the entire system &
            Service failure does not damage the kernel &
            Good \\
        Examples & Linux, BSD & QNX, MINIX & Windows NT, macOS \\
    \end{tblr}
\end{table}

Linux uses a monolithic architecture with support for dynamically loadable modules. This approach gives Linux high performance and flexibility because \term{driver} modules can be added without recompiling the entire kernel.

\begin{figure}[tbp]
    \centering
    \resizebox{0.78\textwidth}{!}{%
    \begin{tikzpicture}
        % Monolithic
        \node[draw, rectangle, minimum width=3.5cm, minimum height=5cm]
            (mono) at (0,0) {};
        \node[above] at (mono.north) {\textbf{Monolithic}};
        \node[draw, fill=red!20, minimum width=3cm, minimum height=4.5cm]
            at (mono.center) {};
        \node[text width=2.5cm, align=center] at (mono.center)
            {Kernel\\[0.5em]All Services};
        \node[below=0.1cm of mono.south, text width=3.5cm, align=center,
            font=\small] {Linux, BSD};
        % Microkernel
        \node[draw, rectangle, minimum width=3.5cm, minimum height=5cm]
            (micro) at (5,0) {};
        \node[above] at (micro.north) {\textbf{Microkernel}};
        \node[draw, fill=red!20, minimum width=3cm, minimum height=1cm]
            at (micro.center) {\small Minimal Kernel};
        \node[draw, fill=blue!20, minimum width=3cm, minimum height=1.2cm]
            at ([yshift=1.5cm]micro.center) {\small Services};
        \node[below=0.1cm of micro.south, text width=3.5cm, align=center,
            font=\small] {MINIX, QNX};
        % Hybrid
        \node[draw, rectangle, minimum width=3.5cm, minimum height=5cm]
            (hybrid) at (10,0) {};
        \node[above] at (hybrid.north) {\textbf{Hybrid}};
        \node[draw, fill=red!20, minimum width=3cm, minimum height=2.5cm]
            at ([yshift=-0.5cm]hybrid.center) {\small Kernel + Some};
        \node[draw, fill=blue!20, minimum width=3cm, minimum height=1.2cm]
            at ([yshift=1.3cm]hybrid.center) {\small Services};
        \node[below=0.1cm of hybrid.south, text width=3.5cm, align=center,
            font=\small] {Windows, macOS};
    \end{tikzpicture}%
    }
    \caption{Kernel architecture comparison: the core block shows services in
    \term{kernel space}, while the separate block at the top shows services in
    \term{user space}}
    \label{fig:kernel-comparison}
\end{figure}

\section{Operating System Evolution}
\label{sec:os-history}

Operating systems evolved alongside hardware progress and user needs. Understanding the broad history explains why modern operating systems are designed as they are today. The main milestones are summarized in \cref{tab:os-history}.

\begin{table}[tbp]
    \centering
    \caption{Broad evolution of operating systems}
    \label{tab:os-history}
    \begin{tblr}{
        colspec = {lX[l]},
        width = \linewidth,
        hlines,
        vlines,
        row{1} = {font=\bfseries},
        cells = {font=\small}
    }
        Era & Main Development \\
        1950s: Early Processing &
            Large computers without an OS; programs ran one at a time using punch cards, without direct user interaction \\
        1969: Birth of Unix &
            Ken Thompson and Dennis Ritchie at Bell Labs created Unix: simple, portable, and written in C with the \textit{do one thing well} philosophy \\
        1980s: PC Era &
            MS-DOS (1981) for IBM PC; Apple Macintosh (1984) introduced a graphical interface; Windows 1.0 (1985) appeared as a layer above DOS \\
        1991: Linux &
            Linus Torvalds from Finland released Linux kernel 0.01 as free and open software; compatible with Unix \\
        1993: Windows NT &
            Microsoft released Windows NT with a hybrid kernel architecture for the corporate and enterprise market \\
        2007--2008: Mobile Era &
            iOS (2007) and Android (2008) redefined computing with touchscreen and mobile-first approaches \\
        2010s: Cloud Era &
            Linux dominated cloud computing; Docker (2013) revolutionized how applications are packaged and distributed \\
    \end{tblr}
\end{table}

Unix became the foundation of almost all modern operating systems. Linux inherits the Unix philosophy, macOS is built on a Unix-based kernel, and Android uses the Linux kernel. The famous 1992 debate between Andrew Tanenbaum, the creator of MINIX and supporter of microkernels, and Linus Torvalds, supporter of monolithic design, showed that both approaches have their place. Linux dominates servers with a modular monolithic design, while QNX with a microkernel dominates automotive systems.

\section{Why Linux?}
\label{sec:why-linux}

Linux is chosen as the main platform in this book for several practical and industry-relevant reasons:

\begin{enumerate}
    \item \textbf{Free and open.} Linux source code is available to everyone. There is no license fee. Every part of the system can be studied directly from its own code.

    \item \textbf{Dominant in servers and cloud.} More than 90 percent of web servers on the internet run Linux. Major cloud platforms such as AWS, Google Cloud, and Azure use Linux as a primary operating system.

    \item \textbf{Transferable concepts.} Linux skills apply to all Unix-based systems, including macOS and BSD. Mastering Linux means having a strong foundation for learning other operating systems.

    \item \textbf{Transparency.} Linux does not hide how it works. Configuration files can be read directly, system logs are available, and every process can be observed with standard commands.

    \item \textbf{Career value.} Demand for workers with Linux expertise continues to grow, especially in DevOps, system administration, and cloud engineering.
\end{enumerate}

This book uses \textbf{Ubuntu Server 24.04 LTS} as the selected Linux distribution. Ubuntu Server is known for stability, complete documentation, and five years of security updates. Ubuntu is also one of the most widely used distributions in commercial cloud environments.

\begin{notebox}
Linux is not one single operating system. Linux is the kernel name. Distributions such as Ubuntu, Debian, Fedora, and Arch Linux all use the same Linux kernel, but differ in \term{package manager}, default configuration, and target users.
\end{notebox}

\section{Installing Ubuntu Server in VirtualBox}
\label{sec:installation}

The practice in this chapter uses a \textbf{\term{virtual machine}}. A \term{virtual machine} is a computer simulated inside a physical computer using software. This approach is safe for learning: if a configuration mistake happens, the \term{virtual machine} can be restored to its original state without affecting the main system at all.

\begin{notebox}
All practice in this chapter is performed from the Ubuntu Server text interface, the command line. There is no graphical interface. This is normal for production servers.
\end{notebox}

\textbf{Prerequisites}: Before starting the installation, make sure the following are available:

\begin{itemize}
    \item VirtualBox version 6.1 or later, installed on the host computer
    \item Ubuntu Server 24.04 LTS ISO file, downloaded from the official Ubuntu website
    \item Host computer with at least 4 GB RAM and 30 GB free disk space
\end{itemize}

\begin{tipbox}
Ubuntu LTS (Long Term Support) means Canonical provides security updates and bug fixes for five years after release. For servers, always use an LTS version so the system stays stable long term.
\end{tipbox}

\subsection*{Lab 1.1 Creating a Virtual Machine}

Create a new \term{virtual machine} in VirtualBox to install Ubuntu Server.

\begin{enumerate}
    \item Open VirtualBox, then click the \textbf{New} button at the top of the window.
    \item Fill in the configuration fields:
    \begin{itemize}
        \item Name: \texttt{Ubuntu-Server-Lab}
        \item Type: \texttt{Linux}
        \item Version: \texttt{Ubuntu (64-bit)}
    \end{itemize}
    \item Set memory (RAM) to \textbf{2048 MB} (2 GB).
    \item Create a new virtual hard disk with size \textbf{25 GB}, VDI type, and dynamic allocation.
    \item After the \term{virtual machine} is created, open \textbf{Settings} and navigate to \textbf{Storage}.
    \item Under Controller: IDE, add an optical drive and choose the downloaded Ubuntu Server ISO file.
    \item Make sure the \textbf{Network} tab uses \textbf{NAT} mode so the \term{virtual machine} can access the internet.
    \item Click \textbf{Start} to power on the \term{virtual machine}.

    \textit{Observe:} The virtual machine boots from the Ubuntu Server ISO and displays the installation menu. This process is identical to starting a physical computer from bootable installation media.
\end{enumerate}

\subsection*{Lab 1.2 Ubuntu Server Installation Process}

Follow these steps after the installation menu opens:

\begin{enumerate}
    \item Choose the language: \textbf{English}.
    \item Choose the keyboard layout that matches the keyboard you use.
    \item On the network configuration screen, keep automatic DHCP settings.
    \item Skip proxy configuration, leaving the field empty.
    \item In the partitioning section, choose \textbf{Use an entire disk}, then confirm to continue. This setting creates partitions automatically.
    \item Fill in the user profile as in this example:
    \begin{lstlisting}[caption={Example user profile fields}]
Your name: Budi Santoso
Server name: ubuntu-server
Username: budi
Password: [use a strong password]
    \end{lstlisting}
    \item Enable \textbf{Install OpenSSH server} so the system can be accessed remotely through a terminal.
    \item Skip additional snap selections and wait for the installation to finish.
    \item When the \textit{Installation complete} message appears, choose \textbf{Reboot Now}.

    \textit{Observe:} The installation process takes about 15 to 20 minutes, depending on internet speed. VirtualBox automatically removes the ISO from the virtual drive when the system restarts.
\end{enumerate}

\subsection*{Lab 1.3 First Login and System Verification}

After the system restarts, perform the first login and verify that the installation succeeded.

\begin{enumerate}
    \item Enter the username and password created during installation.
    \item Run the following commands to verify system information:

    \begin{lstlisting}[language=bash, caption={Verifying system information}]
# Display the running kernel version
uname -r
# Display Linux distribution information
cat /etc/os-release
# Display the IP address obtained from DHCP
ip addr show
    \end{lstlisting}

    \textit{Observe:} The \cmd{uname -r} command displays the active Linux kernel version number. The \cmd{cat /etc/os-release} command displays distribution information. Make sure the distribution name shows Ubuntu 24.04 LTS.
\end{enumerate}

\begin{tipbox}
After installation succeeds, immediately create a snapshot in VirtualBox. Shut down the \term{virtual machine}, open the \textbf{Snapshots} menu, click \textbf{Take}, then name it \textit{Fresh Install}. This snapshot allows recovery to a clean installation state whenever needed without reinstalling from the beginning.
\end{tipbox}

\section{Initial System Configuration}
\label{sec:post-installation}

After the system is installed, several initial configuration steps should be performed before the system is ready for later practice.

\subsection*{Lab 1.4 Updating the System}

System updates ensure that all installed packages are on the latest versions with available security patches.

\begin{lstlisting}[language=bash, caption={Updating system packages}]
# Download latest package information from repositories
sudo apt update
# Upgrade all packages to the latest versions
sudo apt upgrade -y
# Remove packages that are no longer needed
sudo apt autoremove -y
\end{lstlisting}

\textit{Observe:} The \cmd{apt update} command does not install anything. It only updates the list of available packages from \term{repository} servers. The \cmd{apt upgrade} command performs installation based on the updated list.

\subsection*{Lab 1.5 Installing Basic Utilities}

Install several utilities that will be used often throughout the labs in this book:

\begin{lstlisting}[language=bash, caption={Installing basic utilities}]
sudo apt install -y \
    curl \
    wget \
    vim \
    htop \
    tree \
    net-tools \
    neofetch
\end{lstlisting}

After installation finishes, run \cmd{neofetch} to view a summary of system information:

\begin{lstlisting}[language=bash, caption={Displaying system information}]
neofetch
\end{lstlisting}

\textit{Observe:} Neofetch displays the Linux distribution logo together with OS version, kernel, CPU, RAM, and disk information in one compact view. Notice the memory column: how much RAM is already used by the base Ubuntu Server system without additional applications.

\subsection*{Lab 1.6 Verifying Connection and System Resources}

\begin{lstlisting}[language=bash, caption={Checking network and resources}]
# Test internet connection by sending 4 packets to Google DNS
ping -c 4 8.8.8.8
# Display disk usage in human-readable format
df -h
# Display memory and swap usage
free -h
# Display the process list interactively
htop
\end{lstlisting}

\textit{Observe:} The \cmd{ping} command sends four data packets to address 8.8.8.8, Google's DNS server. If responses return with measured \textit{time} values, the internet connection works properly. The \cmd{df -h} command displays each partition's usage in easy-to-read GB or MB units. Press \keystroke{q} to exit \cmd{htop}.

\section{Disk Partitions and Filesystems}
\label{sec:disk-partitioning}

A disk can be divided into several parts called \textbf{partitions}. Think of a filing cabinet with separate drawers: one drawer for work documents, one for personal documents, and one for old archives. Partitions work the same way. Each partition is independent, can use a different \term{filesystem}, and can be protected separately.

\subsection{Partition Tables: MBR and GPT}

Two partition table standards are used today. Their differences are summarized in \cref{tab:mbr-gpt}.

\begin{table}[tbp]
    \centering
    \caption{Comparison of MBR and GPT}
    \label{tab:mbr-gpt}
    \begin{tblr}{
        colspec = {lX[l]X[l]},
        width = \linewidth,
        hlines,
        vlines,
        row{1} = {font=\bfseries},
        cells = {font=\small}
    }
        Aspect & MBR (Legacy) & GPT (Modern) \\
        Firmware & BIOS & UEFI \\
        Disk size limit & 2 TB & 9.4 ZB, practically unlimited \\
        Number of partitions & Maximum 4 primary partitions & Up to 128 partitions \\
        Usage & Older systems & All modern production systems \\
    \end{tblr}
\end{table}

Ubuntu Server 24.04 LTS installation on modern systems uses GPT automatically.

\subsection{Linux Partition Layout}

When selecting \textit{Use an entire disk} in Lab 1.2, Ubuntu automatically creates the following \term{partition layout}:

\begin{description}[noitemsep,topsep=2pt,parsep=0pt]
    \item[\texttt{/boot/efi}] A small EFI partition (512 MB) for storing the GRUB bootloader used by UEFI firmware when the computer starts.
    \item[\texttt{/}] The root partition that contains the entire operating system, programs, and configuration. This is the main Linux partition.
    \item[swap] A disk area used as an extension of virtual RAM. When RAM is full, inactive data is moved to swap to provide space for active processes.
\end{description}

\subsection{Filesystem Types}

Each partition is formatted with a specific \term{filesystem} that determines how data is stored and organized. Common examples in Linux installations are summarized in \cref{tab:filesystems}.

\begin{table}[tbp]
    \centering
    \caption{Common filesystem types in Linux}
    \label{tab:filesystems}
    \begin{tblr}{
        colspec = {lX[l]l},
        width = \linewidth,
        hlines,
        vlines,
        row{1} = {font=\bfseries},
        cells = {font=\small}
    }
        Filesystem & Characteristics & Used for \\
        ext4 & Stable, mature, supports journaling & Root partition (\texttt{/}) and data \\
        swap & Virtual memory area on disk & Swap partition \\
        FAT32 / vfat & Compatible with all firmware & \texttt{/boot/efi} partition \\
    \end{tblr}
\end{table}

\subsection*{Lab 1.7 Checking Partitions and Filesystems}

After the system is installed, inspect the automatically created \term{partition layout}.

\begin{lstlisting}[language=bash, caption={Checking partitions and filesystems}]
# Display disk layout and partitions as a tree
lsblk
# Display partition details including size and type
sudo fdisk -l
# Display filesystem type for each partition
lsblk -f
\end{lstlisting}

\textit{Observe:} The \cmd{lsblk} command displays all disks and partitions in a tree structure. Notice device names such as \texttt{sda1}, \texttt{sda2}, and so on, along with their sizes. The \cmd{lsblk -f} command adds an FSTYPE column that shows the \term{filesystem} for each partition, such as \texttt{ext4} for the root partition and \texttt{vfat} for the EFI partition.

\section{Summary}
\label{sec:rangkuman}

In this chapter, you learned:

\begin{itemize}[noitemsep,topsep=2pt]
    \item \textbf{Operating system functions and types}: an operating system manages processes, memory, files, I/O, and security; there are five OS categories by platform, from desktop and server to mobile, real-time, and embedded.
    \item \textbf{Kernel architecture}: the kernel runs in \term{kernel space} with full privileges; applications in \term{user space} communicate through system calls; the three main architectures are monolithic (Linux), microkernel (QNX), and hybrid (Windows, macOS).
    \item \textbf{Operating system evolution}: Unix (1969) became the foundation of modern operating systems; Linux (1991) inherits the Unix philosophy and now dominates servers and global cloud infrastructure.
    \item \textbf{Ubuntu Server installation}: a \term{virtual machine} in VirtualBox provides a safe learning environment; installation includes VM creation, OS installation, and system verification through \cmd{uname -r} and \cmd{ip addr show}.
    \item \textbf{Partitions and filesystems}: disks are divided into independent partitions; GPT replaces MBR on modern systems; Linux uses ext4 for the root partition, FAT32 for EFI, and swap as \term{virtual memory}.
    \item \textbf{Initial configuration}: system updates with \cmd{apt update} and \cmd{apt upgrade}, basic utility installation, and partition verification with \cmd{lsblk} ensure the system is ready for later labs.
\end{itemize}

\section{Lab Assignment}
\label{sec:tugas}

\textbf{General Instructions}: Complete all exercises independently. Record important steps, save screenshots when needed, and summarize the results as personal documentation.

\subsubsection*{Lab Assignment 1.1 System Installation and Verification}
Install Ubuntu Server 24.04 LTS in VirtualBox following Lab 1.1 through 1.3. After installation finishes:
\begin{enumerate}
    \item Run \cmd{uname -r} and record the displayed kernel version.
    \item Run \cmd{cat /etc/os-release} and record the distribution name and version.
    \item Run \cmd{ip addr show} and record the obtained IP address.
    \item Take a screenshot of each command output above.
\end{enumerate}


\subsubsection*{Lab Assignment 1.2 Initial Configuration and Verification}
After the system is installed, perform the following initial configuration steps:
\begin{enumerate}
    \item Update the system with \cmd{sudo apt update} followed by \cmd{sudo apt upgrade -y}.
    \item Install the \cmd{neofetch} package and run it. Take a screenshot of the result.
    \item Run \cmd{df -h} and \cmd{free -h}. Record the disk and memory capacity available on the system.
    \item Create a VirtualBox snapshot named \textit{Fresh Install} after all steps are complete.
\end{enumerate}


\subsubsection*{Lab Assignment 1.3 Operating System Concept Analysis}
Answer the following questions based on the material in this chapter:
\begin{enumerate}
    \item A company needs operating systems for three purposes: a web server, a medical device with strict response-time requirements, and employee laptops. Which OS category best fits each purpose? Explain why.
    \item Explain the difference between \term{kernel space} and \term{user space}. Why are applications not allowed to access hardware directly?
    \item Unix was born in 1969 and Linux in 1991. Explain the relationship between them and why Unix philosophy remains relevant in today's cloud computing era.
    \item Run \cmd{lsblk -f} on the installed Ubuntu Server system. Record each partition name, \term{filesystem} type, and size. Identify which partition is root (\texttt{/}), which is the EFI partition, and which is swap. Include the command output as a screenshot.
\end{enumerate}
